% !TeX root = ../../Book1.tex
% 纲要标题：### 18.1 Artin 独立性
% 纲要位置：计划纲要.md，第 616 行。

\section{Artin 独立性}
\label{b1:sec:1801}

% 纲要内容（仅作写作计划，尚非教材正文）：
%
% * 乘法特征的线性无关性
% * 域嵌入的线性无关性
% * 与有限群表示特征标正交性的区别
% * 先以最短非平凡线性关系证明任意群上不同乘法特征的有限线性独立性，再得到不同域嵌入的线性独立性；本节不依赖 Galois 对应，且区分函数空间中的线性关系与复特征标内积
%

两个嵌入即使有相同的定义域和值域，也可能作为函数完全不同。我们需要的结论比“两个映射不相等”强：不同嵌入之间不存在系数恒定的非平凡线性关系。证明只依赖乘法，因此先在任意群上建立它。

\subsection{乘法特征的线性无关性}
\label{b1:subsec:1801:01}

设 \(G\) 为群，\(\Omega\) 为域。保持乘法的映射 \(\chi:G\rightarrow\Omega^\times\) 称为一个 \(\Omega\) 值的\term[chengfa-tezheng]{乘法特征}[multiplicative character]。把这些映射视为全体函数组成的 \(\Omega\) 向量空间中的元素，才能谈论线性无关；关系 \(\sum_i c_i\chi_i=0\) 要求同一组常数 \(c_i\) 对每个 \(g\in G\) 都使函数值之和为零。

\ProseParagraphBreak
以下结果称为\term[Artin-dulixing]{Artin 独立性定理}[Artin independence theorem]。
\begin{theorem}[Artin]\label{b1:thm:artin-independence}
设 \(G\) 为群，\(\Omega\) 为域。从 \(G\) 到 \(\Omega^\times\) 的任意一族互异群同态，在全体函数 \(G\rightarrow\Omega\) 组成的 \(\Omega\) 向量空间中线性无关。对无限族，这表示每个有限子族都线性无关。
\end{theorem}
\begin{proof}
若有非平凡关系，从全部这种关系中取非零项数最少者，重新编号写为
\[
c_1\chi_1(g)+\cdots+c_r\chi_r(g)=0\qquad(g\in G),
\]
其中每个 \(c_i\ne0\)。因为 \(\chi_1(g)\ne0\)，不可能 \(r=1\)。由 \(\chi_1\ne\chi_r\)，存在 \(h\in G\) 使 \(\chi_1(h)\ne\chi_r(h)\)。将关系中的 \(g\) 换成 \(hg\)，再减去原关系的 \(\chi_r(h)\) 倍，得到
\[
\sum_{i=1}^{r-1}c_i\Bigl(\chi_i(h)-\chi_r(h)\Bigr)\chi_i(g)=0
\qquad(g\in G).
\]
第一项系数非零，而项数至多为 \(r-1\)，与最小性矛盾。全部操作只涉及有限个项，所以同样排除了无限族中的有限非平凡关系。
\end{proof}
群 \(G\) 不必有限，也不必交换；域的特征没有限制。证明中没有对群元素求平均，也没有除以群阶。

\subsection{域嵌入的线性无关性}
\label{b1:subsec:1801:02}

\begin{corollary}\label{b1:cor:embedding-independence}
设 \(L,\Omega\) 为域。若 \(\sigma_1,\ldots,\sigma_r:L\rightarrow\Omega\) 为互异的域嵌入，则它们在全体函数 \(L\rightarrow\Omega\) 组成的 \(\Omega\) 向量空间中线性无关。
\end{corollary}
\begin{proof}
将嵌入限制到 \(L^\times\)，得到乘法特征。不同域嵌入必在某个非零元素上不同，因为它们都把零映为零。因此这些限制互异，\cref{b1:thm:artin-independence}排除了其上的线性关系，也就排除了原映射的线性关系。
\end{proof}
例如复数上的恒等映射与共轭映射线性无关：若 \(az+b\bar z=0\) 对所有 \(z\in\mathbb C\) 成立，分别代入 \(z=1,i\)，就有 \(a+b=0,a-b=0\)，从而 \(a=b=0\)。一般域上未必有这样显眼的一对检验点，独立性定理保证仍能找到足够多的点。

\subsection{乘法特征与表示特征标}
\label{b1:subsec:1801:03}

本节的乘法特征给出一维线性表示。一般有限群表示的特征标，是给每个群元素的表示矩阵取迹后得到的函数；维数大于一时，这样的迹函数通常不保持乘法。第五卷将研究复数域上不可约表示特征标的正交性，届时需要一个在整个有限群上求和的内积。

\ProseParagraphBreak
这里没有引入那种内积，也没有把“内积为零”和“函数线性无关”混为同一陈述。即使 \(G\) 有限而 \(\Omega\) 的特征整除 \(|G|\)，以 \(|G|\) 作分母的平均运算无意义，本节结论仍然成立。例如特征 \(p\) 的域中，循环 \(p\) 群到乘法群的特征只能是平凡特征，因为 \(x^p-1=(x-1)^p\)；结论并不凭空制造 \(p\) 个不同特征。

\subsection{最短线性关系的证明方法}
\label{b1:subsec:1801:04}

最短关系法先把一个未知的关系压缩到不能再短，再借助结构运算消去一项。关键在于消去某项的同时，必须明确指出另一个系数没有消失；否则新关系可能只是恒等的零关系。下面把函数的独立性转成后文可以使用的矩阵。

\begin{lemma}\label{b1:lem:independent-functions-evaluation}
设 \(X\) 为集合，\(\Omega\) 为域，函数 \(f_1,\ldots,f_r:X\rightarrow\Omega\) 在全体函数 \(X\rightarrow\Omega\) 组成的 \(\Omega\) 向量空间中线性无关。则存在 \(x_1,\ldots,x_r\in X\)，使矩阵 \(\Bigl(f_i(x_j)\Bigr)_{i,j}\) 可逆。
\end{lemma}
\begin{proof}
考虑列向量 \(v(x)=\Bigl(f_1(x),\ldots,f_r(x)\Bigr)^{\mathsf T}\) 在 \(\Omega^r\) 中的张成空间。若其维数小于 \(r\)，对该空间的一组基解齐次线性方程，存在非零行向量 \((c_1,\ldots,c_r)\) 在整个空间上取零值。于是 \(\sum_i c_if_i(x)=0\) 对所有 \(x\) 成立，违背独立性。因此这些列向量张成 \(\Omega^r\)，可以从中选出 \(r\) 个组成一组基，所对应的求值矩阵可逆。
\end{proof}
这里的系数属于值域 \(\Omega\)，不必属于嵌入所固定的较小基域。这一点将在不动域定理的次数估计中起作用。
